\chapter{Preliminaries}

\section{Partial Derivative Refresher}

\subsection{Formal Definition}

The partial derivative of a function $F(x_1, x_2, \ldots)$ with respect to $x_1$ is defined by holding all other arguments fixed:
\begin{equation}
    \pdv{F}{x_1} = \lim_{\Delta x_1 \to 0} \frac{F(x_1 + \Delta x_1,\, x_2, \ldots) - F(x_1, x_2, \ldots)}{\Delta x_1}.
\end{equation}

\subsection{Relations among First Partials}

Consider a function $x(y, z)$. Its total differential is
\begin{equation}
    dx = \left.\pdv{x}{y}\right|_z dy + \left.\pdv{x}{z}\right|_y dz.
\end{equation}

When $z$ is held fixed, this reduces to a one-variable relationship, and the usual inverse-function theorem applies:
\begin{formula}
\begin{equation}
    \left.\pdv{x}{y}\right|_z = \left(\left.\pdv{y}{x}\right|_z\right)^{-1}.
\end{equation}
\end{formula}

A less obvious identity arises when $x$ itself is held constant. Setting $dx = 0$ gives
\begin{equation}
    0 = \left.\pdv{x}{y}\right|_z dy + \left.\pdv{x}{z}\right|_y dz,
\end{equation}
so that $dz/dy = -(\partial x/\partial y)_z / (\partial x/\partial z)_y$. Since $x$ is constant throughout, this ratio is precisely the partial derivative $(\partial z/\partial y)_x$, yielding the \emph{triple product rule}:
\begin{formula}
\begin{equation}
    \left.\pdv{x}{y}\right|_z \left.\pdv{y}{z}\right|_x \left.\pdv{z}{x}\right|_y = -1.
\end{equation}
\end{formula}
The minus sign is a recurring source of error in thermodynamics — it is worth memorising.

\subsection{Mixed Second Partials}

For $x(y, z)$, write $(\partial x/\partial y)_z = f_1(y,z)$ and $(\partial x/\partial z)_y = f_2(y,z)$. Provided $x$ is twice continuously differentiable, the mixed second partials commute:
\begin{equation}
    \left.\pdv{f_1}{z}\right|_y = \left.\pdv{f_2}{y}\right|_z = \frac{\partial^2 x}{\partial y\,\partial z}.
\end{equation}
This symmetry is the key to recognising exact differentials.

\subsection{Exact and Inexact Differentials}

Suppose we are given an expression $df = a(y,z)\,dy + b(y,z)\,dz$ and asked whether a function $f(y,z)$ exists whose total differential it is. A necessary and sufficient condition is that the mixed partials agree:
\begin{equation}
    \left.\pdv{a}{z}\right|_y = \left.\pdv{b}{y}\right|_z.
\end{equation}
When this holds, $df$ is called an \emph{exact differential}.

\subsubsection{Constructing $f(y,z)$ from an Exact Differential}

If the exactness condition is satisfied, $f$ can be recovered by the following procedure. First, integrate $a(y,z)$ with respect to $y$, keeping $z$ fixed and allowing an arbitrary function $g(z)$ to account for the undetermined $z$-dependence:
\begin{equation}
    f(y,z) = \int^y a(y', z)\,dy' + g(z).
\end{equation}
Next, differentiate this expression with respect to $z$ and equate to $b(y,z)$:
\begin{equation}
    b(y,z) = \left.\pdv{f}{z}\right|_y = \int^y \left.\pdv{a(y',z)}{z}\right|_{y'} dy' + g'(z),
\end{equation}
which determines $g'(z)$. Integrating $g'(z)$ then gives $g(z)$ and completes the construction.

\section{Inexact Differentials}

When the exactness condition fails, $df$ does not correspond to any single-valued function of state; rather, its integral depends on both the endpoints \emph{and} the path connecting them. Given points $P_0 = (y_0, z_0)$ and $P = (y, z)$ and a path $C$ between them, the path-dependent quantity is
\begin{equation}
    f_C(P;\,P_0) = \int_C df = \int_C \bigl(a(y',z')\,dy' + b(y',z')\,dz'\bigr).
\end{equation}

It is helpful to think of this as a line integral of the vector field $\vec{V}(y,z) = \bigl(a(y,z),\, b(y,z)\bigr)^T$:
\begin{equation}
    f_C(P;\,P_0) = \int_C d\vec{x}\cdot\vec{V}.
\end{equation}

The path-dependence is most clearly seen by forming a closed loop out of two different paths $C$ and $C'$ between the same endpoints:
\begin{equation}
    f_C - f_{C'} = \oint d\vec{x}\cdot\vec{V}.
\end{equation}
By Stokes' theorem this equals a surface integral of the curl,
\begin{equation}
    f_C - f_{C'} = \int_A d\vec{A}\cdot(\nabla\times\vec{V}),
\end{equation}
where
\begin{equation}
    \nabla\times\vec{V} = \left.\pdv{a}{z}\right|_y - \left.\pdv{b}{y}\right|_z.
\end{equation}
This vanishes — making the integral path-independent — if and only if $df$ is exact. Inexact differentials in thermodynamics (heat $\dbar Q$ and work $\dbar W$) are path-dependent in precisely this sense: the total heat or work exchanged depends on \emph{how} a process is carried out, not merely on the initial and final states.

\chapter{The Dirac Delta Function}

\section{Definition and Basic Properties}

\begin{center}[DIAGRAM: narrow spike of unit area centred at $x = x_0$, representing $\delta(x - x_0)$]\end{center}

\begin{definition}[Dirac delta function]
The Dirac delta $\delta(x - a)$ is defined by the two conditions
\begin{equation}
    \delta(x - a) = 0 \quad \text{for } x \neq a,
\end{equation}
\begin{equation}
    \int_{x_0}^{x_1} \delta(x - a)\,dx =
    \begin{cases} 1 & \text{if } a \in (x_0, x_1), \\ 0 & \text{otherwise.} \end{cases}
\end{equation}
It is not a function in the classical sense but a \emph{distribution}: it is always understood inside an integral sign.
\end{definition}

The defining property is its action on a smooth test function $f(x)$:
\begin{formula}
\begin{equation}
    \int_{x_0}^{x_1} f(x)\,\delta(x - a)\,dx = f(a), \qquad a \in (x_0, x_1).
\end{equation}
\end{formula}
Intuitively, $\delta(x-a)$ ``picks out'' the value of $f$ at $x = a$.

\subsection{Composite Arguments}

The delta function's behaviour under a change of variable is important in practice. For a linear argument $cx$, substituting $u = cx$ shows that
\begin{equation}
    \int f(x)\,\delta(cx)\,dx = \frac{1}{|c|}\,f(0).
\end{equation}
More generally, for $\delta(cx - a)$, the zero occurs at $x = a/c$, giving
\begin{equation}
    \int f(x)\,\delta(cx - a)\,dx = \frac{1}{|c|}\,f\!\left(\frac{a}{c}\right).
\end{equation}

For a nonlinear argument $g(x)$, the delta function has a contribution from each simple zero $x_j$ of $g$ (where $g(x_j) = 0$ and $g'(x_j) \neq 0$), with the local derivative supplying the Jacobian:
\begin{formula}
\begin{equation}
    \int f(x)\,\delta(g(x))\,dx = \sum_j \frac{f(x_j)}{|g'(x_j)|}.
\end{equation}
\end{formula}
This generalisation is used heavily when transforming probability distributions.

\subsection{Indefinite Integral: the Heaviside Step Function}

\begin{center}[DIAGRAM: left panel shows the step function $\theta(x)$ (zero for $x<0$, one for $x>0$) alongside a smoothed approximation; right panel shows the corresponding spike representing $\delta(x)$]\end{center}

Integrating the delta function from $-\infty$ to $x$ produces the Heaviside step function $\theta(x)$:
\begin{equation}
    \theta(x) \equiv \int_{-\infty}^{x} \delta(x')\,dx' =
    \begin{cases} 1 & x > 0, \\ 0 & x < 0. \end{cases}
\end{equation}
Equivalently, $\dv{\theta}{x} = \delta(x)$, so the delta function is the distributional derivative of the step function.

\subsection{Representations}

Because $\delta(x)$ is not an ordinary function, it is often convenient to work with a family of smooth functions that approach it in a limiting sense. The \emph{Gaussian representation} is
\begin{equation}
    \delta_\varepsilon(x) = \frac{1}{\sqrt{2\pi}\,\varepsilon}\,e^{-x^2/2\varepsilon^2}.
\end{equation}
This has unit area for any $\varepsilon > 0$ and width of order $\sqrt{\varepsilon}$, so as $\varepsilon \to 0$ it becomes infinitely narrow while preserving its area, approaching $\delta(x)$ in the distributional sense.

\section{Delta Functions and Functions of a Random Variable}

A powerful application is the transformation of probability distributions. Suppose $x$ is a continuous random variable with probability density $p(x) = dP/dx$, meaning the probability that $x$ lies in the interval $[x, x+dx]$ is $p(x)\,dx$. Given a function $y = g(x)$, we want the probability density $dP/dy$ for the derived random variable $y$.

The derivation uses the delta function as a ``constraint filter.'' We begin from the trivial identity
\begin{equation}
    p(x) = \int dx'\,\delta(x - x')\,p(x'),
\end{equation}
which simply recovers $p(x)$ by the sifting property. The density for $y$ is obtained by replacing the constraint $x = x'$ with $g(x') = y$:
\begin{equation}
    \dv{P}{y} = \int dx'\,\delta\bigl(y - g(x')\bigr)\,p(x').
\end{equation}
Applying the formula for $\delta(g(x))$, let $x_j(y)$ denote all values of $x'$ satisfying $g(x_j) = y$. Then
\begin{formula}
\begin{equation}
    \dv{P}{y} = \sum_j \frac{p(x_j(y))}{|g'(x_j(y))|}.
\end{equation}
\end{formula}
Each pre-image $x_j$ contributes separately, weighted inversely by the local ``stretching'' $|g'|$.

\textbf{Example.} A particle moves at constant speed on a unit circle in the $xy$-plane and is equally likely to be found at any angle $\theta$, so $p(\theta) = 1/2\pi$ for $\theta \in [0, 2\pi)$. We wish to find the distribution of the $x$-coordinate, $x = \cos\theta$.

The density is
\begin{equation}
    \dv{P}{x} = \int_0^{2\pi} d\theta'\,\frac{1}{2\pi}\,\delta(x - \cos\theta').
\end{equation}
The equation $g(\theta') = x - \cos\theta' = 0$ has no solutions for $|x| > 1$, as expected. For $|x| < 1$ there are exactly two solutions: $\theta_1 = \cos^{-1}(x)$ and $\theta_2 = 2\pi - \theta_1$. At both, $|g'(\theta_j)| = |\sin\theta_j| = \sqrt{1-x^2}$. Adding the two contributions,
\begin{equation}
    \dv{P}{x} = 2 \cdot \frac{1}{2\pi} \cdot \frac{1}{\sqrt{1-x^2}} = \frac{1}{\pi\sqrt{1-x^2}}, \qquad |x| < 1.
\end{equation}
The distribution diverges as $x \to \pm 1$ — the particle spends more time near the turning points of its $x$-projection — but the integral over $[-1,1]$ is finite and equal to one.

\chapter{Energy in Thermodynamics}

\section{Thermal Equilibrium}

\begin{definition}[Temperature]
Temperature is the quantity that is equal for two objects that have been in thermal contact for sufficiently long (i.e.\ that have reached \emph{thermal equilibrium}). It can also be understood as the tendency of an object to give up heat to its surroundings: a hotter object transfers energy to a cooler one until equilibrium is reached.
\end{definition}

\begin{definition}[Relaxation time]
The relaxation time is the characteristic time required for two systems in contact to reach thermal equilibrium. It varies enormously between systems — from femtoseconds in dense plasmas to geological timescales in some solids.
\end{definition}

\section{Kinetic Derivation of the Ideal Gas Law}

We derive the ideal gas law from first principles by computing the pressure exerted by $N$ point particles bouncing elastically inside a rectangular box of volume $V$.

\begin{center}[DIAGRAM: single particle in a rectangular box; velocity components $v_x$ and $v_y$ labeled; particle shown reflecting off the right wall]\end{center}

Consider first a single particle moving with velocity $v_x$ toward the right wall. Upon elastic reflection, its $x$-momentum changes by $\Delta p = -mv_x - mv_x = -2mv_x$, so the force on that wall during a single collision is $F_x = -2mv_x/\Delta t$ (averaged over the collision time $\Delta t$).

\begin{center}[DIAGRAM: rectangular box with many particles; shaded slab of depth $d = v_x\Delta t$ adjacent to the right wall of area $A$]\end{center}

With $N$ particles in the box, only those within a distance $d = v_x\Delta t$ of the wall will reach it in time $\Delta t$. This slab has volume $V_{\text{slab}} = v_x\Delta t\cdot A$. Assuming a uniform particle density and that the mean velocity of any large subsection is zero (so half the particles in the slab move toward the wall), the number $N_\partial$ of particles hitting the wall in $\Delta t$ satisfies
\begin{equation}
    N_\partial = \frac{1}{2}\cdot\frac{N}{V}\cdot v_x\Delta t\cdot A.
\end{equation}
The total force on the wall is then
\begin{equation}
    F_x = N_\partial\cdot\frac{-2mv_x}{\Delta t} = -nmAv_x^2,
\end{equation}
where $n = N/V$ is the number density. The pressure $P = |F_x|/A = nmv_x^2$. Because the gas is isotropic, $\langle v_x^2\rangle = \langle v_y^2\rangle = \langle v_z^2\rangle = \tfrac{1}{3}\langle v^2\rangle$, so
\begin{equation}
    P = nm\cdot\tfrac{1}{3}\langle v^2\rangle = \frac{N}{V}\cdot\frac{2}{3}\left\langle\tfrac{1}{2}mv^2\right\rangle = \frac{N}{V}\cdot\frac{2}{3}\,\overline{KE}.
\end{equation}
Identifying $\overline{KE} = \tfrac{3}{2}kT$ (from the equipartition theorem below) yields the ideal gas law:
\begin{formula}
\begin{equation}
    PV = NkT \qquad \text{or equivalently} \qquad PV = nRT,
\end{equation}
where $n$ is the number of moles, $R = 8.31\ \mathrm{J\,mol^{-1}\,K^{-1}}$ is the gas constant, and $k = R/N_A$ is Boltzmann's constant.
\end{formula}

\section{Equipartition of Energy}

\begin{theorem}[Equipartition Theorem]
At temperature $T$, the ensemble average of any energy term that is \emph{quadratic} in a degree of freedom equals $\tfrac{1}{2}kT$.
\end{theorem}

For a monatomic ideal gas the three translational degrees of freedom give $\langle\tfrac{1}{2}mv_x^2\rangle = \langle\tfrac{1}{2}mv_y^2\rangle = \langle\tfrac{1}{2}mv_z^2\rangle = \tfrac{1}{2}kT$, summing to the familiar $\overline{KE} = \tfrac{3}{2}kT$. More generally, each quadratic mode — translational, rotational, or vibrational — contributes $\tfrac{1}{2}kT$ to the average energy. A diatomic molecule, for instance, has two rotational degrees of freedom ($\tfrac{1}{2}I\omega_x^2$ and $\tfrac{1}{2}I\omega_y^2$, with rotation along the bond axis ``frozen out'' by quantum mechanics at room temperature) as well as a vibrational mode that contributes both kinetic and potential energy terms.

\begin{center}[DIAGRAM: left, diatomic dumbbell with two rotational arrows perpendicular to the bond axis; right, the same molecule modeled as two masses on a spring with potential energy terms $\tfrac{1}{2}k_s x^2$, $\tfrac{1}{2}k_s y^2$, $\tfrac{1}{2}k_s z^2$]\end{center}

If $f$ denotes the total number of active quadratic degrees of freedom, the thermal energy of $N$ particles is:
\begin{formula}
\begin{equation}
    U_{\text{thermal}} = Nf \cdot \frac{1}{2}kT.
\end{equation}
\end{formula}

In practice it is safest to use equipartition to compute \emph{changes} in energy $\Delta U$ rather than absolute values of $U$, since zero-point energies and non-quadratic contributions may shift the baseline.
